%%%%%%%%%%%%%%%%%%%%%%%%%%%%%%%
\chapter{Redesign Advection--Diffusion--Reaction Solver}
\label{chap:redesign-adr}

This chapter collects the parts of \code{ADRSolverRedesign} that concern the
\code{SolverCore} framework rather than its discretisation. The unsteady
equation systems derive from \code{SolverCore::UnsteadySystem}; the steady
ones derive from \code{SolverCore::EquationSystem}.

%%%%%%%%%%%%%%%%%%%%%%%%%%%%%%%
\section{Filters}
\label{sec:redesign-adr-filters}

The solver keeps the default field registry of \code{EquationSystem}
(Chapter~\ref{chap:solvercore}): \code{m\_fields}, the transported scalars,
are the solution role and the one physical group, and the history and average
files contain the transported variables. No velocity or pressure role is
published.

The ADR solution represents transported quantities, not fluid momentum. Its
prescribed advection field is an operator coefficient and must not
automatically be advertised as the velocity whose kinetic energy is measured.
Checkpoint, AverageFields, HistoryPoints and Mean are therefore meaningful
through the default registry. If a future coupled ADR system has genuine fluid
state, that equation system can explicitly publish velocity, pressure and
stress roles; the generic scalar ADR solver should not invent them.

SteadyADR derives directly from \code{EquationSystem}, while the filter
lifecycle is driven by \code{UnsteadySystem}, so a steady solve runs no
filters. Supporting them requires moving the lifecycle driver to a common
equation-system or driver hook and defining whether the single solved state is
an initial, step-zero or final event. This is independent of field-registry
compatibility.
